% ====================================================================
% DAFTAR ISI, DAFTAR TABEL, DAFTAR GAMBAR
% ====================================================================
% File ini berisi perintah untuk menggenerate daftar isi, tabel, dan
% gambar secara otomatis dari dokumen.
% LaTeX akan membuat daftar ini berdasarkan \chapter, \section dll
% dan captions gambar/tabel.
%
% CATATAN: File ini TIDAK perlu diedit kecuali ada kustomisasi khusus.
% ====================================================================

% --- DAFTAR ISI (TABLE OF CONTENTS) ---
% Judul otomatis "Daftar Isi", nomor halaman otomatis
\begin{spacing}{1}
\tableofcontents
\end{spacing}
\newpage

% --- DAFTAR TABEL (LIST OF TABLES) ---
% Hanya tampil jika ada tabel dalam dokumen (\begin{table})
% Jika tidak ada tabel, comment baris berikut:
\begin{spacing}{1}
\listoftables
\addcontentsline{toc}{chapter}{DAFTAR TABEL}
\end{spacing}
\newpage

% --- DAFTAR GAMBAR (LIST OF FIGURES) ---
% Hanya tampil jika ada gambar dalam dokumen (\begin{figure})
% Jika tidak ada gambar, comment baris berikut:
\begin{spacing}{1}
\listoffigures
\addcontentsline{toc}{chapter}{DAFTAR GAMBAR}
\end{spacing}
\newpage

% ====================================================================
% CATATAN PENTING:
%
% 1. DAFTAR INI DIBUAT OTOMATIS:
%    - Daftar Isi dari \chapter, \section, \subsection
%    - Daftar Tabel dari \caption dalam \begin{table}
%    - Daftar Gambar dari \caption dalam \begin{figure}
%
% 2. JIKA TIDAK ADA TABEL/GAMBAR:
%    Comment (gunakan %) kedua baris berikut:
%    \listoftables
%    \addcontentsline{toc}{chapter}{DAFTAR TABEL}
%    dan
%    \listoffigures
%    \addcontentsline{toc}{chapter}{DAFTAR GAMBAR}
%
% 3. CARA MENGGUNAKAN DI DOKUMEN:
%
%    Untuk Tabel:
%    \begin{table}[h]
%    \centering
%    \caption{Deskripsi tabel}
%    \label{tab:label}
%    \begin{tabular}{|c|c|}
%    ...
%    \end{tabular}
%    \end{table}
%
%    Untuk Gambar:
%    \begin{figure}[h]
%    \centering
%    \includegraphics[width=0.8\textwidth]{path/to/image}
%    \caption{Deskripsi gambar}
%    \label{fig:label}
%    \end{figure}
%
% 4. COMPILE MULTIPLE TIMES:
%    LaTeX perlu di-compile 2-3 kali agar daftar ini akurat
% ====================================================================
